\subsection{Why a transition norm is a new readout law}
For a closed seed $P_c\psi=\psi$, the closure commutator satisfies
\begin{equation}
 [\widetilde R,P_c]\psi=(I-P_c)\widetilde R\psi=D\psi.
\end{equation}
The right-hand side lies in the orthogonal complementary space. Hence
$\langle\psi,[\widetilde R,P_c]\psi\rangle=0$, even if its
norm is nonzero. Replacing a commutator expectation by
$\|D\psi\|^2$ therefore changes the observable. It is a legitimate
candidate only when that change is stated and physically justified.

The positive matrix $Q=S^*D^*DS$ measures a tendency to leave the
selected closure space under the supplied operator. It does not
assert that an evolving physical state actually leaves that space.
The operator might be a diagnostic rather than the time generator.
The physical mass law and the preservation of the admitted state
domain must be supplied independently.

The distinction between closure and HC parity gives another necessary
test. Under $[P_c,J]=0$, define
\begin{equation}
 R_o=\tfrac12(\widetilde R-J\widetilde RJ),\qquad
 D_o=\tfrac12(D-JDJ)=(I-P_c)R_oP_c.
\end{equation}
Only $D_o$ changes HC grading. A pacer commuting with $J$ has
$R_o=0$ and cannot supply this channel, even if it has nonzero
closure defect. A selfadjoint comparison coupling can be formed as
$A_o=D_o+D_o^*$, but forming it is an additional operation; its
appearance in the native equations is not established by its definition.

\subsection{Closure penalties do not determine the missing channel}
Let $M$ be a bounded represented matrix and split the carrier into
the range and kernel of $P_c$. Write
\begin{equation}
 M=\begin{pmatrix}A&F\\G&H\end{pmatrix},\quad
 K_c=[M,P_c]=\begin{pmatrix}0&-F\\G&0\end{pmatrix},\quad
 E_c=\begin{pmatrix}0&F\\G&0\end{pmatrix}.
\end{equation}
Direct block multiplication gives
\begin{equation}
 K_c^*K_c=E_c^*E_c=
 \begin{pmatrix}G^*G&0\\0&F^*F\end{pmatrix}.
\end{equation}
Thus a minimal positive penalty
$H_c=\rho K_c^*K_c+\sigma E_c^*E_c$ is block diagonal and
commutes with $P_c$. Its compression can measure the norm of a
transition generated by $M$, but the penalty itself has no off-diagonal
closure channel. Positivity of the penalty cannot replace the missing
represented pacer block.

If the pacer were a norm limit of polynomials in an exactly closed
$M,M^*$, then it would commute with $P_c$ as well and its transition
would vanish. Conversely, exact closure of $M$ says nothing about
an independently specified pacer. One needs a bridge between those
operators before transferring a closure statement from one to the other.

There is a particularly sharp freedom result even when the pacer is
unitary. Given any positive contraction $0\leq Q\leq I_3$, define
\begin{equation}
 R_Q=\begin{pmatrix}\sqrt{I_3-Q}&-\sqrt Q\\
                         \sqrt Q&\sqrt{I_3-Q}\end{pmatrix},\qquad
 P_c=\operatorname{diag}(I_3,0).
\end{equation}
The square roots commute because both are functions of $Q$.
Block multiplication gives $R_Q^*R_Q=I_6$, and the transition from
the first three coordinates to the complement has Gram matrix $Q$.
Therefore unitarity, positive readout and the dimension three leave
every positive contraction spectrum available. These assumptions
do not select a hierarchy or the Koide radius.

\subsection{When a Gram law follows from a gapped block}
A possible physical route can be expressed more concretely at a
frozen background. Assume a canonically normalized Hermitian mass
operator on a retained three-dimensional space and its complement:
\begin{equation}
 \Omega_{\rm tr}=\begin{pmatrix}A&V^*\\V&H_\perp\end{pmatrix},
 \qquad A=A^*,\quad H_\perp=H_\perp^*\geq\Delta I,
 \quad \Delta>0.
\end{equation}
For an eigenvalue $|\lambda|<\Delta$, solve the lower equation to obtain
\begin{equation}
 y=-(H_\perp-\lambda I)^{-1}Vx,
\end{equation}
and substitute into the upper equation:
\begin{equation}
 [A-V^*(H_\perp-\lambda I)^{-1}V]x=\lambda x.
\end{equation}
This is an energy-dependent Schur complement. The resolvent identity
controls the error made by replacing it by its static value:
\begin{equation}
 \|A_{\rm eff}(\lambda)-A_{\rm eff}(0)\|
 \leq\frac{\|V\|^2|\lambda|}{\Delta(\Delta-|\lambda|)}.
\end{equation}
The bound uses the two inverse norms
$\|(H_\perp-\lambda I)^{-1}\|\leq(\Delta-|\lambda|)^{-1}$
and $\|H_\perp^{-1}\|\leq\Delta^{-1}$. It is a controlled
comparison only in the stated gap regime.

With the further assumptions $A=0$, $H_\perp=\Delta I$ and
$V=\mu B$ with a real coupling scale, each eigenvalue $q_i$ of $Q=B^*B$ gives a light
signed eigenvalue $-m_i$, where
\begin{equation}
 m_i=\frac{\sqrt{\Delta^2+4\mu^2q_i}-\Delta}{2}
 =\frac{\mu^2}{\Delta}q_i
  +O\left(\frac{\mu^4q_i^2}{\Delta^3}\right).
\end{equation}
Thus a Gram mass-magnitude law can be a leading approximation in
a canonical Hermitian Dirac-type reading. The signed effective
operator is negative in this comparison; its interpretation as a
positive physical mass magnitude is an additional reading, not
a sign that can be dropped without comment.

The exact eigenstate is dressed. Its complementary norm satisfies
\begin{equation}
 \frac{\|y\|^2}{\|x\|^2}=\frac{m_i}{\Delta+m_i},\qquad
 \|y\|^2=\frac{m_i}{\Delta+2m_i}
\end{equation}
after normalization. Enforcing strict closure $y=0$ would instead
require $Vx=0$. A nonzero light mass produced in this zero-direct-block
mechanism is consequently incompatible with exact closure in the
original retained projector. Approximate closure or a separately
defined dressed spectral projector are possible alternatives, but
they are different state constraints.

\subsection{Cyclic compatibility fixes a form, not its coefficients}
Suppose a unitary carrier operation $W$ and the generation three-cycle
$U$ satisfy
\begin{equation}
 WS=SU,\qquad WP_c=P_cW,\qquad W\widetilde R=\widetilde RW.
\end{equation}
Then $BU=WB$ and hence $QU=UQ$. For a nonzero transition Gram matrix, the positive root $L=\sqrt Q$
also commutes with $U$. In three dimensions the three distinct Fourier
eigenvalues of $U$ imply that every such Hermitian root has the form
\begin{equation}
 L=aI_3+zU+\overline zU^*,\qquad a>0,\quad z=re^{i\theta}.
\end{equation}
Its eigenvalues are
\begin{equation}
 \ell_k=a+2r\cos(\theta+2\pi k/3).
\end{equation}
Trace orthogonality gives
\begin{equation}
 \operatorname{Tr}L=3a,\qquad
 \operatorname{Tr}L^2=3a^2+6r^2,\qquad
 K=\frac13+\frac{2r^2}{3a^2}.
\end{equation}
The Koide value therefore requires $r=a/\sqrt2$. Cyclic compatibility
does not supply this coefficient ratio. Conversely, once the native
$Q$ were represented, its coefficients would be determined without a
mass fit by
\begin{equation}
 a=\tfrac13\operatorname{Tr}\sqrt Q,\qquad
 z=\tfrac13\operatorname{Tr}(U^*\sqrt Q).
\end{equation}
The native burden is thus a joint operator calculation, not the
selection of a familiar trigonometric parametrization.

At the Koide radius positivity gives three phase chambers,
\begin{equation}
 \theta\in\bigcup_{j=0}^2
 [2\pi j/3-\pi/12,\ 2\pi j/3+\pi/12]
 \pmod{2\pi}.
\end{equation}
They follow by intersecting the inequalities
$\cos(\theta+2\pi k/3)\geq-1/\sqrt2$ for all three $k$.
At the boundary $\theta_b=3\pi/4$, write
$\theta=\theta_b-\delta$ inside the chamber. The small root is
$\ell_0=a\delta+O(\delta^2)$, so its candidate mass is
$m_0=\kappa a^2\delta^2+O(\delta^3)$. A light eigenmass can
therefore arise near a boundary in this comparison. Neither the
distance to the boundary nor its dynamical preservation has yet been
selected by the native theory.

\subsection{Deriving the positivity obstruction in the quadratic model}
For $\eta\geq0$, the quadratic root model can be written as a linear function of the
two-mode density matrix $\rho=\chi\chi^*$. In the Fourier basis of
$U$, its signed root eigenvalues are
\begin{equation}
 \ell_k=\frac{X_0}{\sqrt3}
 +\sqrt{\frac23}\eta[X_1\cos\varphi_k-X_2\sin\varphi_k],
 \qquad \varphi_k=2\pi k/3.
\end{equation}
Equivalently $\ell_k=\sqrt3\operatorname{Tr}(F_k\rho)$, where
\begin{equation}
 F_k=\frac{I_2+\sqrt2\eta(\cos\varphi_k\sigma_z-
                               \sin\varphi_k\sigma_x)}{3}.
\end{equation}
The three $F_k$ sum to $I_2$. The Pauli combination has eigenvalues
$\pm1$, so $F_k$ has eigenvalues $(1\pm\sqrt2\eta)/3$.
Positivity for every $\rho\geq0$ is therefore equivalent to
$\eta\leq1/\sqrt2$. Necessity follows by testing a negative
eigenvector if this condition fails; sufficiency follows because
the trace of a product of two positive matrices is nonnegative.
This is the precise sense in which a globally positive linear-density
readout in this family cannot achieve the intended isometric radius.

Replacing the input metric by the pullback metric of this map would
make a chosen nonzero $\eta$ an isometry by definition. It would
not derive its physical value. Nor does covariance under the three-cycle
fix the relative singlet and doublet weights: the constant direction
and the zero-sum plane are separate invariant pieces. A comparison
tensor can assign them different positive weights while preserving
the cyclic symmetry.

The antipodal $5/9$ bound is similarly a property of the actual
root eigenvalues. If $cI_3\pm D$ are both positive, the three
eigenvalues $d_i$ obey $|d_i|\leq c$ and $\sum_i d_i=0$.
The maximum of $\sum_i d_i^2$ on this planar cube section is
$2c^2$, attained at permutations of $(c,0,-c)$. Consequently
\begin{equation}
 \frac{\operatorname{Tr}[(cI_3\pm D)^2]}
      {[\operatorname{Tr}(cI_3\pm D)]^2}
 =\frac{3c^2+\operatorname{Tr}D^2}{9c^2}\leq\frac59.
\end{equation}
Selecting one allowed positive branch can evade the pairwise
requirement. A physical explanation then has to supply that branch
selection rather than infer it from the involution alone.

\subsection{Phase overlap: root averaging and mass averaging differ}
The published overlap comparison provides another explicit test of
what is being averaged. In one common generation basis, let
\begin{equation}
 x_{jk}=a[1+\sqrt2\operatorname{Re}(\omega^ke^{i\delta_j})],
 \quad \omega=e^{2\pi i/3},\quad
 w_j\geq0,\quad\sum_jw_j=1.
\end{equation}
Each preparation is assumed to have nonnegative roots and the same
mean $a$. Averaging the roots before squaring gives
\begin{equation}
 \overline x_k=a[1+\sqrt2\operatorname{Re}(\omega^k Z)],
 \qquad Z=\sum_jw_je^{i\delta_j},\qquad
 K=\frac{1+|Z|^2}{3}.
\end{equation}
For two equally weighted phases $\delta_0\pm h$, this reduces the
radius by $\cos h$ when $|h|<\pi/2$, leaving the central phase.
It follows from a vector average on the unit circle, not from a
new dynamical mass equation.

Averaging squared roots gives a different answer:
\begin{equation}
 \sum_jw_jx_{jk}^2=\overline x_k^2+
                      \sum_jw_j(x_{jk}-\overline x_k)^2.
\end{equation}
The variance term cannot be removed by calling both procedures a
superposition. A coherent coupling of two paths is yet another
specified construction. For $B(\phi)=B_1+e^{i\phi}B_2$,
\begin{equation}
 B(\phi)^*B(\phi)=B_1^*B_1+B_2^*B_2+
 e^{i\phi}B_1^*B_2+e^{-i\phi}B_2^*B_1.
\end{equation}
Here relative phases can affect the operator spectrum through cross
terms. A common phase multiplying all of $B$ cancels from its Gram
matrix. Changing only the prepared probe can change spectral weights
without moving the operator poles. These alternatives must be
distinguished before a transport interpretation is assigned.

An inverse reconstruction can find phase mixtures consistent with
specified masses and an assumed intermediate gap. Its feasibility
region is useful for excluding combinations of assumptions, but the
reconstructed phases and gap are not independent predictions. The
native operator, the admitted preparations and the rule for combining
paths must be fixed before the resulting spectrum is compared with
data as a test of the model.

\subsection{A bent octant and a flux-free core (exploratory reading)}

The norm relation of the main text fixes the Koide value but not the
three states that carry it. Take three unit states with pairwise
overlap $|\langle a|b\rangle|^2=\tfrac12$ and write
$G_{12}G_{23}G_{31}=e^{i\varphi}/(2\sqrt2)$ for their Bargmann
invariant. Reading the roots as eigenvalues of the Gram matrix $G$,
\begin{equation}
 K=\frac{\operatorname{Tr}G^2}{(\operatorname{Tr}G)^2}=\frac23
 \quad\text{for every }\varphi,\qquad
 \det G=\tfrac12\bigl(\sqrt2\cos\varphi-1\bigr).
\end{equation}
At $\varphi=\pi/4$ the determinant vanishes: the three states lie in
one two-dimensional space, one root is zero, and the configuration is
the octant triple of three spin-$\tfrac12$ states along orthogonal
axes. For $\varphi<\pi/4$ the states span $\mathbb C^3$ and all three
roots are positive. The value $\varphi=\tfrac23$, which is the phase
obtained from the charged-lepton pole masses, reproduces the root
ratios $1:14.38:58.97$. The overlap $\tfrac12$ and the phase are inputs
here; nothing in this paragraph derives them.

Geometrically, on a Bloch sphere three points at mutual angle $90^\circ$
form a rigid octant whose enclosed solid angle $\pi/2$ fixes
$\varphi=\pi/4$. A phase below $\pi/4$ with the same overlaps requires
leaving the sphere, that is, a genuinely three-dimensional state
space. The deficit $\pi/4-\varphi$ then acts like a flux removed from
the octant.

The long Companion offers one carrier for such a deficit: a guarded
core excised from the Hopf fibration, $X=S^3\setminus B_\varepsilon(K)$,
which projects to a cap on the Bloch sphere. Read as a cap of solid
angle $\Omega$ that contains no flux, three statements follow for a
spin-$\tfrac12$ Berry connection. First, with the cap inside the
triangle the phase falls below $\pi/4$; outside it does not. Second,
on a sphere without $n$ caps there are $n-1$ free Aharonov--Bohm
holonomies, so exactly one core is the only number that bends the
phase without introducing a free parameter. Third, the displaced flux
stays on the sphere; spread evenly it does not give the lepton value,
whereas placed outside the triangle it gives
$\varphi=\pi/4-\Omega/2$.

The Companion calls the operation that sends a state to a definite
partner state the \emph{reflector}, or Umkehrwalze. The name is taken
from the Enigma cipher machine, which is a cryptographic device and not
a physical model: its reflector returns every incoming signal at a
different contact, never at the same one, so it is an involution
without fixed points. Only this mathematical property is used. The
Companion identifies the physical counterpart as an antilinear
involution containing time reversal. On a spin-$\tfrac12$ sphere an
antilinear map $\Theta$ with $\Theta^2=-1$ is the antipodal map, the
only such map without fixed points, whereas $\Theta^2=+1$ gives a
reflection with fixed points. Read in this way, the core returns the
displaced flux at its antipode, point to point, which is the case
$\varphi=\pi/4-\Omega/2$.

The status of this reading should be stated plainly. The required
deficit $\pi/4-\tfrac23=0.1187$ lies numerically close to half the
horizon counting parameter $\gamma\approx0.2375$ used in the Companion,
but this was found with the lepton masses known, and the size of the
core is not derived. It coincides with the open question why the
counting rule differs between the matter and the tachyon reading.
Core size and the fraction of displaced flux placed outside the
triangle are degenerate; a value $\Omega=\gamma$ holds only together
with complete placement outside. The construction is therefore a
geometric reformulation of an input, not a mass prediction.

\subsection{Where the missing number comes from (status after further tests)}

After the first edition of this reading, six pre-registered tests
tried to obtain the deficit from the carrier itself: a singlet
admixture, nested normalisations, entangled generations,
generation-dependent frames, an observer algebra and a time octant
built from the frame maps. All were refuted. The reason is a rigidity:
symmetric constructions fix the Bargmann phase by their overlap
(spin $\tfrac12$, $1$, $\tfrac32$, $2$ give $0.785$, $0.679$, $0.650$,
$0.637$), and $\tfrac23$ is not a rational part of $2\pi$. A flux-free
core escapes this rigidity because it changes the connection on the
generation sphere, not the states. Topologically, one core is the only
rigid bend, and an antipodal reflector makes it a single puncture in
$\mathbb{RP}^2$; its size, however, cannot be fixed topologically. It
has to be supplied from outside.

In the Companion that outside number is the horizon counting
parameter of the matter reading, $\gamma=0.237533$, obtained by
counting every spin $j$ with weight two in
$\sum_j g_j\,e^{-2\pi\gamma\sqrt{j(j+1)}}=1$. This rule is chosen, not
derived. Counting all $2j+1$ states instead gives $0.274067$, and the
conditional tau mass of the main text would move to about
$1868$\,MeV. Replacing $\sqrt{j(j+1)}$ by $j$ or $j+\tfrac12$ changes
the weight-two value to $0.3497$ or $0.2206$. The tau-mass test is
therefore a joint test of the core reading, the weight-two counting
and the area spectrum. Two partial results bear on the counting: the
weight two is compatible with full rotation invariance if the readout
averages over the axis, and a register model in which every strand of
a puncture carries the state of its frame reproduces the full count
$2j+1$ for a maximally mixed reading but not the constant two. Why the
matter reading counts with weight two remains open.

\subsection{Orientation of the octant step}

The bent-octant reading uses one Bargmann phase, $\pi/4$, for the flat
octant. Five later tests, registered in advance but with outcomes fixed
analytically, asked what selects this phase in the carrier. For the
three spin-$\tfrac12$ rays along $+x$, $+y$, $+z$ the ordered overlap
product is $(1+i)/4$. The reverse loop $x\to z\to y\to x$ has the same
endpoints and the same three transition probabilities $\tfrac12$, but
the conjugate product. Magnitudes and endpoints therefore fix the size
of the step, $\pi/4$, but not its sign. Constant step sequences
$\pm1$ modulo eight visit all eight phases in one cycle, so a uniform
distribution of accumulated phases does not show independent steps.

The compensator does not supply the sign either. On the closed odd
generation states $JS=-S$, so $C=(I+J)/2$ annihilates them. A real
projection of rank five that removes one real direction of the complex
null line $L$ of the octant readout exists, $C_n=I-P_n$, and leaves the
readout unchanged; but the two real directions $n$ and $in$ of $L$ give
the same ranks and the same final readout, and $J$ does not choose
between them. At the photon point of the frame convention ($b=0$,
$\rho=I_2/2$) the inputs leave the Bargmann phase free. Only if the
phase is identified with the Bures distance $\tfrac12\arccos r$ from
the pure boundary does the Gram determinant become
$\det G(r)=(\sqrt{1+r}-1)/2$, with rank two exactly at $r=0$; that
identification is assumed, not derived. A colour-equivariant linear
readout of the generation raster into a singlet vanishes, so colour
does not select the phase.

What remains is a selection of the ordered return among admissible
continuations. A two-window model makes the timing explicit: the first
window keeps the admissible orientations as a mask inside one
provisional state, the second applies a second mask, and a sign is
emitted only if exactly one orientation survives. Of the sixteen mask
pairs, three give $+\pi/4$, three give $-\pi/4$, nine are blocked and
one is ambiguous. The existing repair and persistence rules of the
Companion are compatible with this timing but contain no quantity that
decides the sign. The octant phase therefore has an open orientation
in addition to the open size of the core discussed above. A later
check locates this orientation: on the invariant doublet the frame maps
and parity keep it, while time reversal of the field amplitudes acts as
complex conjugation and reverses it. The sign is therefore a time
orientation of the readout, and no structure recorded in the Companion
fixes it.
A later series of explorations, none of them blind, asked whether a further
ingredient fixes this orientation: a zero point of the puncture area, a
residual rotation of a three-arc disc read as a holonomy, light as a
mediator between the frames, a perturbation of the light weight, and a
rotation rate of the disc as the origin of $\Lambda$ (layer v4.10.0-r59).
Each computation assumes the orientation; none fixes it. The open choice
ends at the choice of the imaginary unit, which over $\mathbb C$ is a
convention. A relative orientation of two sectors is not: a test
registered in advance (H125) confirms, as a structure computation with an
expected outcome, that the family $su(2)$ of the sedenion construction is
the conjugate doublet on the leptons relative to the Lorentz spin and not
on the quarks, with Bargmann invariants $\tfrac14\mp\tfrac i4$ in all
admissible assignments. Under the setting that CP violation arises in the
same way from the family structure in both sectors this gives
$\sin\delta_{\rm PMNS}<0$, a prediction blind only against future data
and open until a determination at $3\sigma$ (deadline 2035). Reading the
field value of the successive Hilbert shifts as Cayley--Dickson algebras
gives $U(1)$, $SU(2)$ and $G_2$ with $su(3)$ on the first levels and only
a discrete $S_3$ beyond; it fixes no number.

\subsection{Family symmetry, octonion copies and the Koide point on the tower}

The explorations and tests that followed the handedness series are
collected in layers v4.10.0-r60 to r65. Most parts are not blind; the
layers record the blindness per part, and refuted predictions are
named below as refuted.

On the Albert algebra $H_3(\mathbb O)$, with its three off-diagonal
octonions read as generations (a reading, not derived), the commutant
of $G_2$ in $F_4$ is a family $SO(3)$. The charge per slot breaks it to
exactly $A_4$, and this breaking is a single spin-3 tensor (test H127:
this part confirmed, the part on a leak into the quarks refuted by its
wording). On the colour triplets the $A_4$-invariant and the
$SO(3)$-invariant operators coincide, so with exact $A_4$ the quarks
have degenerate masses and trivial mixing; quark mixing needs $A_4$
itself to break. The vacuum alignment depends on free coefficients of a
potential, so no blind test of it is possible at present
(layer v4.10.0-r60).

One Cayley--Dickson doubling over the octonions contains exactly three
$G_2$-stable octonion copies, conditional on the definition of a copy;
the closure condition is $\operatorname{Im}(t_1+it_2)^3=0$. The counts
predicted in advance for the next two levels are refuted. The three
copies are not three generations: no common charge preserves them, and
together they cover 15 of 24 complex dimensions, so they are three
equivalent readings of one generation. On the next level four of the
nine octonionic copies can leave a complete further generation as a
remainder that cancels in the algebra; the advance prediction that no
such selection exists is refuted, and whether remainder states are
physical states is not shown. The copy structure alone gives no masses
($|K-\tfrac23|\geq0.206$ in the Gram readout), and the real line on
which the remainder is built is a neutral scalar, not a spin-2
graviton. The number of generations remains an input
(layer v4.10.0-r62).

In the amplitude readout $K=1/(3\cos^2\theta)$ for the mixing angle
$\theta$ between the real unit and the copy direction, so the Koide
value is exactly $\theta=\pi/4$. This angle is the fixed point of an
antilinear exchange $J$ of matter-type and tachyon-type states. With a
set scale and with $J$ required as a symmetry, $K=\tfrac23$ holds at
tree level and the running cancels at first order in the logarithm. The algebra
does not supply $J$: it allows unequal mass terms, and $J$ is neither
an automorphism nor an antiautomorphism. Equal mass terms are thus
equivalent to a postulate of equal standing of the two frames, which
the author adopted as a working assumption on 26 November 2026; it is
not derived. Without a set scale $K$ depends on the starting point.
The phase $\delta=2/9$ remains a fit with one parameter. The measured
value is $K=0.6666645\pm0.0000051$, $0.43\sigma$ below $\tfrac23$, with
the error entirely from the tau mass; if $K=\tfrac23$ holds exactly for
pole masses, then $m_\tau=1776.969$\,MeV, a conditional prediction
(layer v4.10.0-r63). The exchange $J$ acts as $z\mapsto i\bar z$ on the
plane concerned; it is neither the tachyon map $z\mapsto iz$ of the
frame convention nor time reversal, and the value $1777.134$\,MeV of
the main text explains a deviation that is not significant at present
(layer v4.10.0-r64).

Layer v4.10.0-r64 also records the tests of the same days that belong
to other lines. For the Hermitian $3\times3$ matrices over the
sedenions the derivation algebra has dimension 22 (numerical rank; the
value 17 fixed in advance is refuted). Mixing built from the copies
under named settings gives one class, a trimaximal column without CP
violation (the predicted CP violation is refuted). The layer records a
rule of the author: before a statement that a prediction is disfavoured
by data, the prediction is to be transferred to the matter frame. With
the frame mask the class is compatible with the measured
$\theta_{12}$ only if the mask also acts on the neutrinos; this is
compatibility, not a prediction. An open prediction is recorded as
assigned: under the setting that every horizon crossing is one quantum,
area steps $\Delta A=4\ln2\,\ell_P^2$ and hence echoes after the
ringdown (H105, to be decided by gravitational-wave analyses, deadline
31 December 2028). The value $4\ln2$ is not specific to the
construction.

The charge axis of the sedenion construction is a placeholder. Two
generations per carrier, the family $su(2)$ and the opposite
orientation of leptons and quarks of the preceding subsection hold for
every axis and for both types of charge, left and right multiplication
(test H162, confirmed, largely not blind; layer v4.10.0-r65). The
choice of the axis stays a setting between two named criteria. The
same layer restricts the statement on the relative orientation: it
does not depend on the choice of the imaginary unit, the axis or the
charge type, but it does depend on the setting that the same side of
the multiplication carries the family in both sectors. That setting is
the natural one and does not follow from charge and colour alone; with
different sides the two sectors are equally oriented. The prediction
$\sin\delta_{\rm PMNS}<0$ therefore rests on this setting as well; its
status and deadline are unchanged.

\subsection{The lepton phase: an angle, a branch and one measured number}

From layer v4.10.0-r69 of the Companion on, the frames of the frame
convention carry names that say what they are: the rest frame R
(formerly reference frame), the matter frame M, the dual frame T
(formerly tachyon frame), which is the image of M under the duality
$z\mapsto iz$, and the light point P (formerly photon eigenframe), a
fixed point of the state disk rather than the frame of an observer.
Under the sector assignment H99, which remains a reading, the two
lepton sectors are called the Re sector (neutral leptons) and the Im
sector (charged leptons); the quarks keep the names up and down sector.
The names are a convention and move nothing. Layer v4.10.0-r70 places
the frames on the sphere of field states of the doublet
$(\mathbf E,\mathbf B)$. The compensator sees only the height
$b_C=\mathbf E\cdot\mathbf B/u$ on its axis, with
$u=\tfrac12(|\mathbf E|^2+|\mathbf B|^2)$: R lies at the pole $b_C=1$,
M at $b_C=1/\gamma_L$, P on the equator and T at $b_C=-1/\gamma_L$. The
same layer corrects a reading of r69: the port weights $1:x$ and an
amplitude ratio $x$, with $x=\beta^2$, describe different states, and the
quasiconformal dilatation $\gamma_L$ belongs to the second, not to the
matter frame. No verdict changes.

Layer v4.10.0-r73 records fifteen tests on the lepton phase, each
registered in advance. In each sector $s$ the masses are written in the
Brannen form
\begin{equation}
 \sqrt{m_k}=\mu_s\bigl(1+\sqrt2\cos(\delta_s+2\pi k/3)\bigr),\qquad k=0,1,2,
\end{equation}
which gives $K=\tfrac23$ for every phase. The form is set, and
$K=\tfrac23$ rests on the postulate of equal standing of the frames. The
phase of the Im sector is $\delta^*=0.2222220$, the value $2/9$ of the
preceding subsection; it lies $x=\pi/12-\delta^*=0.0395773$ below the
octant value $\pi/12$. The neutrino ratio
$R=\Delta m^2_{21}/\Delta m^2_{31}$ gives an offset
$x_\nu=0.04336\pm0.00213$ of the Re phase from $\pi/6$. With pole
masses, the only scheme in which $K=\tfrac23$ holds, the two offsets
differ by $1.78\sigma$; a figure of $4.2\sigma$ in an earlier research
note used the wrong uncertainty and is withdrawn.

In a phase potential
$V_s(\delta)=-\alpha'_s\cos2(\delta-\chi_s-\varphi)+q\cos3\delta$, with
$(\alpha'_s,\chi_s)=(1,\pi/4)$ in the Im sector and $(-1,0)$ in the Re
sector, a common rotation $\chi\to\chi+\varphi$ moves both minima exactly
by $\varphi$, to $\pi/12-\varphi$ and $\pi/6-\varphi$, for every
$\varphi$. A sector scale or a conjugation term does not reproduce both
offsets (misses of $9.2\sigma$ and $23\sigma$). The common offset is
therefore a dimensionless rotation angle, not a scale. Each minimum is
still twofold. The cubic spectral invariant
\begin{equation}
 \prod_{k=0}^{2}\bigl(1+a\cos(\delta+2\pi k/3)\bigr)=\frac{4-3a^2+a^3\cos3\delta}{4}
\end{equation}
with the rule ``larger $\cos3\delta$'' selects the branch of the data in
both sectors with the same sign; the rejected Re branch would miss $R$
by $51\sigma$. The rule is the limit $q\to0^-$ of the sixth-order term:
the scale drops out and only a sign remains. That sign follows neither
from $J$, nor from stability, nor from a natural parent term; the author
set it on 28 November 2026.

The value of $\varphi$ is not derived. It is not a Wigner rotation of
the frame boosts (the nearest of 48 candidates fixed in advance lies
$118\sigma$ off), and none of four simple terms that break the exchange
of the sectors fixes it. Nor is it the minimum of the common vacuum
energy of the two sectors: that energy has the exact exchange symmetry
$V_{\rm Re}(-\delta;\pi/4-\varphi)=V_{\rm Im}(\delta;\varphi)$ and settles
at $\varphi=\pi/8$, where both sectors would have the same spectrum. The
measured value breaks this exchange. Thus $\varphi=0.0395773$ stays the
one measured number of the lepton picture.

With $K=\tfrac23$ and the branch sign set, three inputs, $\varphi$ and
the scales $\mu_e$ and $\mu_\nu$, carry all six lepton masses (test H227,
carries partly). The Koide relation gives $m_\tau=1776.969$\,MeV from
$m_e$ and $m_\mu$ ($+0.43\sigma$), and $\varphi$ gives $R=0.030839$
($+1.80\sigma$); neither is blind. With $\mu_\nu$ fitted to
$\Delta m^2_{31}$ the neutrino masses follow in normal ordering,
$m_{1,2,3}=0.379$, $8.81$ and $50.11$\,meV, with
$\sum m_\nu=59.30\pm0.24$\,meV, $m_\beta=8.87$\,meV and $m_{\beta\beta}$
between $1.29$ and $4.00$\,meV for free Majorana phases. All lie below
the present bounds and are recorded as predictions; a measured
$\sum m_\nu<58$\,meV or an established inverted ordering would refute the
picture. The ratio of the two scales is no integer power of a frame
number.

Layer v4.10.0-r74 records nine further tests registered in advance, H228
to H236, and one exploration. None derives $\varphi$. As the limit of a
loop of frame operations, $\varphi$ is no attracting fixed point (the
nearest lies $27\sigma$ off), and as a stage of the gradient flow towards
$\pi/8$ it is reached at none of ten natural times fixed in advance
(H228). Loosening the postulate $J$ at order six by admitting the full
square of the cubic norm selects the data branch through stability, so
the sign setting could be traded for a changed rule; but the minima of
the vacuum energy then lie only at $\pi/12$, $\pi/8$ and $\pi/6$ (H229).
No declared frame map takes the symmetry point $\pi/8$, which lies in R,
or the point $0$ to the measured value in M: maps that keep the spectral
form keep the phase, and the others destroy $K=\tfrac23$ (H230). A field
$\varphi$ rolling with Hubble friction reaches the measured value today
only for a mass of $0.97\,H_0$, but it would then still be rolling, and
the resulting drift of the electron mass, $2.7\cdot10^{-10}$ per year, is
excluded by atomic clocks by more than six orders of magnitude (H231). A
backward search shows that, for every term that depends only on
$\delta$, $\varphi$ equals the mean of the two sector offsets at every
equilibrium; a symmetric term can make the measured value a local but
never a global minimum (H232). The fine angles do not lock onto a raster
of eight positions with up to eight halvings ($\varphi$ lies at least
$392\sigma$ from every grid point; H233), and $\varphi$, $\beta-\tfrac12$
and two other fine values share no common small step beyond chance
(H234). A Casimir twist energy between an inner and an outer field, with
counting weights two for the Im and one for the Re sector, breaks the
exchange of the sectors but leaves the data point at most a local
minimum (H235). The exploration k449 behind this test, not blind, reads
two separated regions as the split-complex algebra $\mathbb R\oplus\mathbb Rj$
with $j^2=1$, their merger as $\mathbb R$, and a continuous phase across
the gap as $\mathbb C$, which is put in by hand; in the Casimir gap only
the mixed boundary condition is repulsive, with the familiar ratio
$-7/8$. It supplies no number.

The last test of the series, H236 (carries), clears a double counting.
Up to H235 several computations carried both the sixth-order coupling
$q=-0.0852$ fitted in an earlier test and the rotation $\varphi$, each of
which can produce the offset. The two offsets fix the division:
$\varphi$ carries the offset, and $q$ is only the sign of the branch
choice, $q\to0^-$; the data allow $|q|\leq0.029$ at $2\sigma$ but do not
require any. No verdict changes, and the mass balance H227 with its
predictions is not affected. The open choice has thus stayed where layer
r73 left it, at a measured angle and a set sign. OP3 and OP4 remain
open.
